% Pro Quinctio (pro-quinctio) — English translation
% Translated by Claude (Anthropic) from the Latin (Perseus, Clark 1909).
% Style guide: STYLE.md.

\ciceroSection{1} The two things that count for most in our state — the two — are working against us in this case, both of them: standing at the highest pitch, and eloquence. One of these, Gaius Aquilius, troubles me; the other, I genuinely fear. That the eloquence of Quintus Hortensius will hinder me as I speak --- this gives me some unease. That the standing of Sextus Naevius will harm Publius Quinctius --- this terrifies me, and not in moderation.

\ciceroSection{2} Nor would there be such cause for complaint that they enjoy these things in their highest measure, were our own at least middling. As it is, the case stands like this: I, who have neither sufficient experience nor great natural ability, am pitted against the most fluent advocate in Rome; and Publius Quinctius — whose resources are slender, whose means are nothing, whose circle of friends is small — must contend with the most well-connected adversary one could imagine.

\ciceroSection{3} A further disadvantage attaches to us. Marcus Iunius, who has argued this case before you on several occasions before now --- a man both well-practised in other cases and thoroughly versed in this one --- is at this moment away, detained on a fresh embassy. The case has fallen to me; and even granting that I might bring everything else in adequate measure, I have certainly had barely enough time to learn so substantial a matter, tangled in so many controversies.

\ciceroSection{4} The thing that customarily comes to my aid in other cases, then, fails me in this one too. For the deficit of natural ability I have habitually made up by sheer diligence; and how great that diligence is cannot be appreciated unless time and scope are granted to display it. The greater these disadvantages mount, Gaius Aquilius, the more carefully you and your council must hear our words --- so that truth, weakened by many handicaps, may at last be revived by the equity of men such as yourselves.

\ciceroSection{5} For if you, the judge, shall be seen to have offered no shelter to isolation and want against force and influence, if before this council a case is to be weighed by means rather than by truth, then surely there is nothing left holy or pure in our state, nothing that can offer the lowly any consolation in the gravity and virtue of the judge. Surely either with you and these men sitting beside you the truth will prevail, or, driven from this place by force and influence, it will find no resting-place at all. I do not say this, Gaius Aquilius, because your loyalty and constancy are in any doubt with me, or because Publius Quinctius ought not to repose his highest hope in these picked men of the state whom you have summoned. What is the matter, then?

\ciceroSection{6} First: the magnitude of his danger has struck the man with extreme fear, since by a single judgment everything he has is decided. As he turns this over, the thought of your power comes to him no less often than the thought of your fairness; for everyone whose life is placed in another's hands thinks more often about what the man in whose dominion and power he stands \emph{can} do than about what that man \emph{ought} to do.

\ciceroSection{7} Second: Publius Quinctius has Sextus Naevius for an adversary in name only --- in reality he has the most fluent, most forceful, most flourishing men of our generation, who by united effort and with all their resources are defending Sextus Naevius --- if defending is the right word for accommodating one man's appetite so that he can the more easily crush by an unjust verdict whomever he chooses.

\ciceroSection{8} For what more inequitable, what more shameful thing, Gaius Aquilius, can be uttered or even mentioned, than that I, who am to defend another's life, his reputation, his fortunes, must speak first? Especially when Quintus Hortensius, who in this trial holds the place of the prosecutor, will speak against me --- a man on whom nature has lavished the highest abundance and ease of speech. So it has come about that I, whose proper task is to deflect blows and heal wounds, am compelled to perform that task before the adversary has thrown a single weapon; while they are given the time of attack at the very moment when no power is left us either to dodge their assault or, if they shall do what they are prepared to do — fling some false charge as though some envenomed weapon — to apply any remedy.

\ciceroSection{9} This has come about by the praetor's inequity and injustice. First, because contrary to all custom he chose that judgment on the man's character should precede judgment on the substance; second, because he so set up that very judgment that the defendant should be compelled to argue his case before he had heard a word from his accuser. And this was done at the instance of those whose influence and power are such that, just as if their own substance or honour were at stake, they so devotedly accommodate Sextus Naevius's enthusiasm and appetite, and so freely test out their own resources in cases of this kind, in which --- the more they can do because of their virtue and noble birth --- the less they ought to display how much they can do.

\ciceroSection{10} Burdened and battered, then, by so many and such great difficulties, Publius Quinctius has taken refuge, Gaius Aquilius, in your good faith, your truthfulness, your mercy. Up to this point he has not been able to find equality of right against the violence of his adversaries, nor equal capacity to act, nor a magistrate fairly disposed; up to this point everything has been hostile and harmful to him through the highest injustice. He prays you, Gaius Aquilius, and you who sit beside him in council, that you allow equity --- thrown about and harassed by many wrongs --- in this place at last to find a footing and be made firm.

\ciceroSection{11} That you may do this the more easily, I will take pains to make plain to you from the beginning how the affair was conducted and contracted. Gaius Quinctius was the brother of this Publius Quinctius. In the rest of his affairs he was a careful and attentive paterfamilias, in one matter rather less considered: that he formed a partnership with Sextus Naevius. Naevius is a respectable enough man, but not so trained that he could know the rights of partnership and the duties of a settled paterfamilias. Not for any want of natural ability — Sextus Naevius was never reckoned an unwitty buffoon nor a coarse auctioneer. What is the truth, then? When nature had given him nothing better than a voice, and his father had left him nothing but his freedom, he turned his voice to profit and used his freedom for licence to be impudent.

\ciceroSection{12} For which reason, indeed, the only thing you would gain by taking him on as your partner was to put him to school in your money on what the profits of money might be. Yet, drawn in by old habit and intimacy, Quinctius made (as I have said) a partnership in the affairs that were being managed in Gaul. He had a substantial cattle business, and an estate in the country well farmed and very productive. Naevius is taken from the Licinian halls and from the auctioneers' bench, and ferried off to Gaul, all the way across the Alps. The change of place was great; the change of character, none. For a man who from his boyhood had set himself to make money without outlay, once he had laid out something or other and put it to the joint stock, could not be content with a moderate profit.

\ciceroSection{13} Nor is it any wonder, if a man whose voice had been for sale supposed that what he had earned by his voice was destined to bring him a great profit. And so, by Hercules, he was not moderate in funnelling whatever he could from the joint funds into his own private house --- in this matter as scrupulous as if those who conduct partnership in good faith are accustomed to be condemned in arbitration on the partner's behalf. But on these matters I have no need to say what Publius Quinctius is eager to have me mention. The case demands it, but because it merely demands and does not insist, I will pass over it.

\ciceroSection{14} When the partnership had now run for several years, and when Naevius had often been suspect to Quinctius and could not easily render an account of those things he had managed by appetite rather than by reason, Quinctius dies in Gaul, with Naevius present, and dies suddenly. By his will he left as heir this Publius Quinctius, so that the man on whom the highest grief came by his death might also receive the highest honour from the same source.

\ciceroSection{15} On his brother's death, not very long afterwards, Quinctius sets out for Gaul. There he lives on familiar terms with this Naevius. They are together for nearly a year, frequently consulting between themselves about the partnership and about the whole of that Gallic account and business. In all that time Naevius let drop not a single word that the partnership owed him anything or that Quinctius personally had been in his debt. Since some debt had been left behind, on accounts of which money would need to be raised at Rome, Publius Quinctius here gives notice that he intends to hold an auction at Narbo, in Gaul, of those things that were his own private property.

\ciceroSection{16} Then this excellent Sextus Naevius — at much length — dissuades the man from holding the auction. He could not sell so easily at the time he had advertised; the means in money were available to himself at Rome; and these, if Quinctius had any sense, he would consider their joint resource, in view of that brotherly closeness and his own connection by marriage. (Naevius has a cousin of Publius Quinctius's as wife, and children by her.) Because Naevius was saying what an honourable man ought to say, Quinctius believed that one who imitated the speech of honourable men would imitate their actions too. He drops the plan to hold the auction; he sets out for Rome. Naevius, leaving Gaul, comes to Rome at the same time.

\ciceroSection{17} Since Gaius Quinctius had owed money to Publius Scapula, Publius Quinctius settled, with you as arbiter, Gaius Aquilius, what should be paid out to Scapula's children. The settlement was managed through you because, owing to the state of the treasury, it was not enough to have looked at the books to see how much was owed: one also had to ask at the temple of Castor how much was actually being paid. You decide, and on account of your friendship with the Scapulae, you fix exactly what should be paid them, to the last denarius.

\ciceroSection{18} All this Quinctius was doing on Naevius's prompting and advice. Nor is there any wonder if he relied on the counsel of one whose support he believed assured. For Naevius had not only promised in Gaul, he was saying every day at Rome that, the moment Quinctius gave him the nod, he would pay the money down. Quinctius for his part saw that Naevius could do it, understood that he ought to do it, and did not suppose he would lie, since there was no reason for him to lie. As if he had the cash at home, so he set down with the Scapulae what he would give them. He notifies Naevius and asks him to see to what he had said.

\ciceroSection{19} Then this excellent fellow --- I am afraid he will think himself made fun of, that for the second time now I call him excellent --- thinking he had brought his man into the tightest of corners, and meaning at the very crisis of the moment to bind him to his own terms, says he will not give a copper coin until first all the affairs and accounts of the partnership had been settled, and he had ascertained that no controversy with Quinctius was going to remain. ``We will look at those things later, Quinctius,'' he says. ``For the present I would prefer that you yourself see to what you have undertaken, if you don't mind.'' Quinctius says he will not act on any other terms. What he had promised, says Naevius, was no more binding on him than if, when conducting an auction, he had promised something at the owner's bidding.

\ciceroSection{20} Crushed by that desertion, Quinctius gets a few days' grace from the Scapulae and sends to Gaul that the things he had advertised should be sold. He holds the auction in his absence at a worse moment, and pays off the Scapulae on harder terms. Then he himself approaches Naevius and says that, since he suspected there would be some controversy on some point, Naevius should see to it that the whole business was wound up as soon and as painlessly as possible.

\ciceroSection{21} Naevius gives as his friend Marcus Trebellius, our common acquaintance; we put up Sextus Alfenus, our relative --- a man brought up in his house and one whom Naevius habitually used a great deal. The matter could in no way be settled, because our man wanted to take a moderate loss, and Naevius was not content with moderate plunder.

\ciceroSection{22} From that moment forward, then, the affair began to proceed by sureties for appearance. When the dates set for appearance had often been postponed and a fair amount of time had been used up to no purpose, Naevius came in for the next appearance. I beg you, Gaius Aquilius, and you who sit beside him in council, give your closest attention, that you may recognise an unprecedented kind of fraud and a new method of laying snares.

\ciceroSection{23} He says he had held the auction in Gaul; that what had seemed to him to be his, he had sold; that he had taken pains to leave nothing owing to himself from the partnership; that he was no longer requiring any sureties for appearance from Quinctius nor giving any himself; if Quinctius wished to take any action against him, he would not refuse. At this, Quinctius — eager to revisit his Gallic affairs — does not bind the man over for the present; and so they part without sureties for appearance. Then Quinctius stays in Rome for about thirty days; he postpones the other sureties for appearance that he had outstanding, so as to be free to go to Gaul; he sets out. He left Rome two days before the Kalends of February.

\ciceroSection{24} Quinctius did so, that is, on the 29th of January, in the consulship of Scipio and Norbanus. I beg you to commit that day to memory. Lucius Albius, son of Sextus, of the Quirina tribe — a respectable man and of the first rank for honour — set out with him. When they had come to the place called Ad Vada Volaterrana, they meet a close acquaintance of Naevius's, who was bringing slaves for sale from Gaul to him: Lucius Publicius. He, on coming to Rome, tells Naevius where he had seen Quinctius.

\ciceroSection{25} On hearing this from Publicius, Naevius sends his slaves around among his friends. He himself rounds up his own associates from the Licinian halls and the entrance to the meat-market, asking them to be at his side at the second hour next day at the Sextian Bank. They turn up in numbers. Naevius testifies that Publius Quinctius had not appeared and that he himself had appeared. The records are signed and sealed in the chief place by the seals of distinguished men, and they break up. Naevius applies to the praetor Burrienus to be allowed, by his edict, to take possession of Quinctius's goods. Burrienus orders the goods to be advertised --- the goods of one with whom Naevius had had close intimacy, partnership, family connection through children still living, things that could in no way be torn apart.

\ciceroSection{26} From this it could easily be understood that there is no obligation, however sacred and solemn, that greed does not habitually weaken and violate. For if friendship is cultivated by truth, partnership by good faith, kinship by affection, then a man who has tried to strip a friend, a partner, a kinsman of his reputation and his fortunes must necessarily admit that he is empty, faithless, and wicked.

\ciceroSection{27} Sextus Alfenus, agent for Publius Quinctius --- a close associate and relation of Sextus Naevius --- tears down the bills of advertisement, takes back a single slave whom Naevius had seized, and gives notice that he is the agent, that Naevius ought to consider the reputation and fortunes of Publius Quinctius and to wait for his return; that if Naevius will not do so but has settled on coercing him by such tactics into accepting his terms, Alfenus on his side will plead nothing, and if Naevius wishes to take any action, will defend him by trial.

\ciceroSection{28} While these things are being done at Rome, Quinctius in the meantime --- contrary to law, custom, and the praetor's edicts --- is forcibly driven out of their common pasture and field by their common slaves. Suppose, Gaius Aquilius, that everything Naevius has done at Rome was done with restraint and reason: only if what was done in Gaul on his written orders looks rightly and properly done. Driven and ejected from the property after suffering an outstanding wrong, Quinctius takes refuge with Gaius Flaccus, the commander, who at that time was in the province --- a man whom, as his rank requires, I name with honour. How vigorously he thought the matter ought to be redressed you can learn from his decrees.

\ciceroSection{29} Alfenus meanwhile was fighting daily at Rome with that ageing gladiator Naevius, and was using a public substantial enough, since Naevius was unrelenting in his pursuit of Quinctius's life. Naevius demanded that the agent give security for satisfying judgment. Alfenus says it is not fair for an agent to give security in a matter where the principal would not be required to give it if he were present. The tribunes are appealed to. From them aid was obtained for certain, and so the affair is broken off on these terms: that on the Ides of September Sextus Alfenus would undertake to produce Publius Quinctius.

\ciceroSection{30} Quinctius comes to Rome and discharges his bond. This man --- this most ferocious creature, possessor, expeller, plunderer of goods --- for a year and six months claims nothing, lies quiet, leads our man on with negotiations as long as he can. Then at last he applies to the praetor Gnaeus Dolabella that Quinctius be required to give him security for satisfying judgment under the formula: ``in respect of which he sues a man whose goods have been in possession for thirty days under the praetor's edict.'' Quinctius did not refuse to be ordered to give such security, if his goods had been possessed under the edict. Dolabella decrees --- I will say nothing of how fair it was, only this: that it was unprecedented; and even this I would rather have left unsaid, since either point was perceptible to anyone --- that Publius Quinctius should enter into a sponsio with Sextus Naevius: ``if his goods were not in possession for thirty days under the edict of Praetor Burrienus.'' Those who were then attending Quinctius objected, and pointed out that the issue ought to be tried on its merits, so that either both should give security to each other or neither: there was no necessity to bring the reputation of one party alone into the trial.

\ciceroSection{31} Quinctius himself protested loudly that the reason he was unwilling to give security was so as not to seem to have judged that his goods had been possessed under the edict; further, that if he were to enter into the sponsio in such terms, he would --- as is in fact what is happening now --- have to argue a case for his life and reputation in the prior place. Dolabella --- in the manner of noble persons, who, whether they begin to act rightly or wrongly, are so excellent in either course that no one of our station could match them --- proceeds to commit the wrong with the utmost firmness. He orders us either to give security or enter the sponsio, and meanwhile orders our advocates, who were objecting, to be sharply removed.

\ciceroSection{32} Quinctius leaves much shaken, as well he might, who had been given a choice as miserable as it was inequitable: either himself to condemn himself in a capital matter (had he given security), or to argue the case for his life in the prior place (had he made the sponsio). Since in the one course there was no escape from rendering the gravest possible judgment --- one's own judgment of oneself --- but in the other there was hope of coming before such a man as judge as could bring the more help the less he could bring of personal influence, he preferred to make the sponsio. He made it. He chose you, Gaius Aquilius, for judge, and brought the action on the bond. In this the whole essence of the trial and the case rests.

\ciceroSection{33} You see, Gaius Aquilius, that this is a trial not about money, but about the reputation and fortunes of Publius Quinctius. And whereas our forebears so framed things that the man who pleaded for his head should plead in the latter place, you understand that we are pleading in the prior place, and on the basis of an unheard prosecution charge. Further, you see those who have been our defenders turned to prosecuting us, and natural gifts which used to be turned to bringing safety and aid being turned to ruin. One thing only remained --- and they did it yesterday: to take you to the praetor in order to have him fix in advance the time we might be allowed for speaking. They would readily have got that from him, had you not made known what your jurisdiction and your duty and your power are.

\ciceroSection{34} And up to this point we have had no one but you with whom to maintain our right against them, and they have never been content with maintaining what could win general approval. They suppose that influence stripped of injustice is shallow and weak. Since, however, Hortensius is pressing on you to retire to consider, asking that I not consume time with my speech, complaining that under our previous advocate the case could never be brought to a peroration --- I will not allow that suspicion to remain, that we do not wish the case to be judged. I will not claim that I can present the case more conveniently than it has been presented before, nor will I make many words about it. The reason is that the case has been already laid out by the man who spoke then, and brevity is required of me, who can neither think up nor utter a great deal --- a brevity that is, in any case, my closest friend.

\ciceroSection{35} I shall do what I have often noticed you doing, Hortensius. I shall divide the whole presentation of my case into definite parts. You always do this, because you always can; I shall do it in this case, because in this case I think I can. What nature gives you in always being able to do it, the present case grants me in being able to do it today. I shall set myself definite limits and boundaries, beyond which I cannot go even if I most want to, in order that I may have a fixed subject on which to speak, that Hortensius may have the points laid out to which he must reply, and that you, Gaius Aquilius, may be able to know in advance the matters on which you are about to hear evidence.

\ciceroSection{36} We deny, Sextus Naevius, that you took possession of Publius Quinctius's goods under the praetor's edict. The sponsio has been made on that issue. I shall show, first, that there was no cause for which you could apply to the praetor to be allowed to take possession of Publius Quinctius's goods; second, that you could not have taken possession under the edict; and last, that you did not take possession. I beg you, Gaius Aquilius, and you who sit in council, to commit carefully to memory what I have promised. You will the more easily grasp the whole matter if you keep this in mind, and you will easily call me back by your judgment if I try to step beyond the bounds I have set for myself. I deny there was any cause for the application; I deny you could have taken possession under the edict; I deny you took possession. When I have shown these three things, I shall make my peroration.

\ciceroSection{37} There was no cause for the application. How can this be made clear? Because Quinctius owed Sextus Naevius nothing on the partnership account, nor anything privately. Who is the witness to this? The same person who is its bitterest adversary. On this matter I will summon you, you yourself, Naevius, as my witness. For a year and more after Gaius Quinctius's death, Quinctius was in Gaul together with you. Show that you ever asked him for that vast unheard-of sum of money, show that you ever made mention of it, ever said it was owed --- and I will grant that it was owed.

\ciceroSection{38} Gaius Quinctius dies, who, as you say, owed you a great sum on definite accounts. Publius Quinctius, his heir, comes into Gaul, to you in person, to your common land, in fact to the very place where not only the property but every account and every paper was. Who would be so loose in his domestic affairs, who so neglectful, who so unlike you, Sextus, that, when an estate had passed from the man with whom he had contracted to that man's heir, he would not, the first time he saw the heir, inform him, address him, render his account, and, if anything came into dispute, settle it either privately or by strict law? Is that how it is? What good men do, if they wish their kinsmen and connections to be dear to others and to be regarded as honourable --- this Sextus Naevius did not do, who burns and is borne along by greed to such a pitch that he is willing to give up some part of his own advantage so as not to leave any part of any honour to this kinsman of his?

\ciceroSection{39} And he would not, of course, demand money if any were owed --- he, who, because what was never owed has never been paid, is now trying to snatch away not just money but the lifeblood and very life of a kinsman? You did not, of course, wish to bother the man whom now you do not let breathe freely; the man you now wickedly wish to kill, you were too modest to dun. So I am to believe: a kinsman, attentive to you, an honourable man, a respectful man, your senior --- you were unwilling, or didn't dare, to dun him. Often, as happens, having steadied yourself, having resolved to mention the money, having come up prepared and rehearsed, this timid soul with maidenly modesty would suddenly hold himself back; the words would suddenly fail; you would want to address him but didn't dare, lest he hear unwillingly. That, no doubt, is what happened.

\ciceroSection{40} Are we to believe that Sextus Naevius, who is now attacking the man's life, spared his ears in those days? If he had owed it, Sextus, you would have demanded it --- demanded it at once. If not at once, then a little later. If not a little, then somewhat later. Within those six months, surely. By a year's end, beyond doubt. But after a year and six months, when you had daily access to the man to remind him, you say not a word. After almost two full years are gone, you finally accost him. What ruined and prodigal spendthrift, with a fortune not yet eaten away but actually overflowing, was ever so dilatory as Sextus Naevius? When I name the man, I think I have said enough.

\ciceroSection{41} Gaius Quinctius owed you; you never demanded it. He died; the property passed to the heir. Though you saw him daily, only after two years do you finally address him. Will it be doubted whether it is more probable that Sextus Naevius would have demanded what was owed at once, or that he would not have addressed his man even after two years? Was there no time to address him? But he lived with you for over a year. Could the matter not be conducted in Gaul? But justice was administered in the province, and trials were held in Rome. The remaining options are that either the highest negligence stood in your way, or unique generosity. If you say negligence, we shall be amazed; if you say goodness, we shall laugh. I cannot find what else you might say. There is sufficient proof that nothing was owed to Naevius, in the fact that for so long he demanded nothing.

\ciceroSection{42} What if I show that this very thing he is now doing is a witness that nothing is owed? What is Sextus Naevius doing now? What is the controversy about? What is this trial in which we have now been engaged for two years? What business is being conducted in which he is wearing out so many men of such standing? He is asking for money. Now at last? Very well, let him ask; let us hear him.

\ciceroSection{43} ``He wants the accounts and disputes of the partnership decided.'' Late --- but at some point, after all. Granted. ``That,'' he says, ``is not what I am pursuing, Gaius Aquilius, and not what I am working at now. Publius Quinctius has been using my money for so many years.'' Let him use it freely; I am not asking. Why then are you fighting? Is it for what you have often said in many places --- so that he should be no citizen, no longer hold the place he has hitherto held with the highest honour, no longer be counted among the living? So that the verdict on his life and all his honours may be passed; so that he may plead his case before a judge in the prior place; so that, only when he has finished, he may at last hear his accuser's voice? What? What is the point of this? That you may sooner have what is yours? But if you wanted that, the matter could long ago have been concluded. So as to fight a more honourable contest?

\ciceroSection{44} But you cannot cut the throat of Publius Quinctius, your kinsman, without the highest crime. So that the trial may be easier for you? But Gaius Aquilius does not gladly give judgment in a capital matter, and Quintus Hortensius has not learned to speak against a man's life. What, then, is being asked of us, Gaius Aquilius? He demands money; we deny it is owed. Let there be a trial at once: we do not refuse. Anything else? If he is anxious about whether the proceeds of judgment will be ready when judgment is given, let him take security for satisfying judgment; in the same words in which he takes security from me, let him give security on what I demand. The thing can already be over with, Gaius Aquilius; you can already walk free of a burden almost as heavy as Quinctius's.

\ciceroSection{45} Where do we stand, Hortensius? What do we say to this offer? Can we at last lay down our arms and dispute about money without the parties' fortunes being on the line? Can we so pursue our own claim that we leave a kinsman's life intact? Can we take on the part of plaintiff and lay down the part of accuser? ``No,'' he says, ``I will take security from you; but I will not give security to you.'' Who, pray, allots us terms so even as that? Who has decreed this --- that what is fair against Quinctius is unfair against Naevius? ``Quinctius's goods,'' he says, ``were possessed under the praetor's edict.'' So you ask, in order that I might admit it, that what we are arguing in court has never been done, we should --- as if it had been done --- confirm by our own action.

\ciceroSection{46} Can no way be found, Gaius Aquilius, that each may come as soon as possible to what is his own without disgrace, infamy, or destruction to anyone? Surely if anything were owed he would demand it, would prefer some other procedure to this single one out of which all these proceedings are sprung. The man who in all those years did not even address Quinctius, when he could have done so daily; who, when he first did begin to act, used up all his time in postponing the appointed dates; who afterwards even let the bond go; who by ambush forcibly threw him off their common land; who, when he could have brought the matter to trial without any objection, preferred a sponsio about character; who, when he is recalled to the very trial out of which all this has grown, repudiates the fairest of terms --- such a man as good as openly admits that he is seeking not money but life and blood. Such a man openly says: ``If anything were owed me, I would demand it, and would already long ago have got it;

\ciceroSection{47} I would not need so vast an undertaking, so invidious a trial, so abundant a body of advocates, if it were a matter of demanding. It must be wrung out unwillingly; what is not owed must be torn out and squeezed; Publius Quinctius must be tipped off all his fortunes; the powerful, the eloquent, the noble must all be summoned. Force must be brought against truth, threats hurled, dangers held up, terrors interposed, so that, conquered and frightened by these things, he gives himself up at last.'' All which, by Hercules, when I see those who fight against us and reflect on this assembly, seems to be at hand and impending and impossible to escape; but when I have turned my eyes and mind back to you, Gaius Aquilius, the more vehemently and zealously these things are pursued, the lighter and feebler I judge them. Nothing, then, was owed --- as you yourself proclaim.

\ciceroSection{48} What if it had been owed? Would there have been ground for applying to the praetor for permission to take possession of the goods? I do not think so --- in justice, nor to anyone's profit. What does Naevius show, then? He says his bond was broken. Before I prove that this was not done, I would like, Gaius Aquilius, to consider both the affair itself and Sextus Naevius's own action in the light of his duty and of universal practice. ``He had not come to the bond'' --- as you say --- ``the man with whom you had connection by marriage, partnership, every old reason and tie of intimacy''. Was it right at once to rush to the praetor? Was it right at once to apply for leave to take possession of the goods under the edict? Were you so eager to descend to those last and most hostile remedies that you reserved nothing more grave or cruel for what you might do later?

\ciceroSection{49} What more shameful thing can happen to a man, what more bitter and miserable to a man of standing? What disgrace can be greater, what calamity comparable? If fortune has taken money from someone, or someone's injustice has snatched it away, still, while reputation is intact, honour easily consoles want. Even one who is afflicted with disgrace or convicted in a shameful trial enjoys at least his own goods --- though he has, what is most miserable, no expectation of another man's resources --- in his miseries he is yet relieved by this aid and comfort. But the man whose goods have been put up for sale, the man whose ample fortunes — and not only those: even his necessary food and clothing — have been put up under the auctioneer with disgrace, this man is not only thrust out of the number of the living, but, if it can be done, is sent away below the dead too. For an honourable death often even adorns a shameful life; a life so shameful does not leave even an honourable death any room.

\ciceroSection{50} And so, by Hercules: when a man's goods are taken into possession under the edict, his whole fame and reputation is taken with the goods. The man whose advertisements are posted in the most thronged places is not allowed to perish even in silence and obscurity. The man for whom managers are appointed and proprietors set up to declare on what law and on what conditions he is to be sold; the man over whom the auctioneer's voice declaims and fixes a price --- to him, while still alive and seeing, the bitterest funeral is announced; if it is to be called a funeral when not friends gather to honour the obsequies, but buyers of goods, like executioners, gather to tear and rend the remains of his life.

\ciceroSection{51} Therefore our forebears wished this rarely to happen, and the praetors arranged that the thing should be done with great consideration. Honest men, even when they are openly being defrauded, even when there is no power of bringing them to trial, descend to that step timidly and step by step, compelled by force and necessity, unwilling, after many bonds have been broken, often after they themselves have been tricked and let down. For they consider how serious a thing it is, and how great, to advertise another man's goods. To cut a citizen's throat, even by the law, no good man wants. He prefers it remembered that, when he could have ruined a man, he spared him, rather than that, when he could have spared him, he ruined him. These things good men do for the most distant strangers, in the end even for their bitterest enemies, both for the sake of human reputation and for the sake of common humanity, so that, having done no harm to another knowingly, no harm by right may justly fall on themselves.

\ciceroSection{52} ``He did not come to the bond.'' Who? A kinsman. If the matter taken by itself were grave, still its severity would be lightened by the name of kinship. ``He did not come to the bond.'' Who? A partner. You ought to extend something even more handsome to the man whom either choice associated with you, or fortune linked. ``He did not come to the bond.'' Who? The man who had always been at your service. So, against him who only this once committed the error of not being at your service, you have hurled all the weapons that are kept in store for those who have done many wrongs and committed many frauds?

\ciceroSection{53} If a wretched twopence of yours were at stake, Sextus Naevius, if in some trifling matter you feared some entrapment, would you not at once have run to Gaius Aquilius or to one of those gentlemen who give legal opinions? When the rights of friendship, partnership, kinship were at stake, when account had to be taken of duty and of reputation, you not only did not refer to Gaius Aquilius or Lucius Lucilius --- you did not even consult yourself; you did not even say this much to yourself: ``It is two hours past the time. Quinctius has not come to his bond. What am I to do?'' By Hercules, if you had said these two words to yourself --- ``what am I to do?'' --- your appetite and avarice would have caught its breath, you would have given some little space for reason and consideration; you would have collected yourself; you would not have come to such a nasty pass as to have to confess before such men as these that, in the same hour in which the bond fell to you broken, you took the resolve of overturning a kinsman's fortunes from the foundations.

\ciceroSection{54} I now consult these gentlemen on your behalf, after the time and on someone else's matter, since you in your own matter forgot to consult, when there was time. I ask you, Gaius Aquilius, Lucius Lucilius, Publius Quinctilius, Marcus Marcellus: ``A certain partner and kinsman of mine, with whom I have an old close tie and a recent dispute over money, has not come to his bond. Do I apply to the praetor to be allowed to take possession of his goods --- or, when his house, wife, and children are at Rome, do I rather give notice at the house?'' What is the answer that any of you could possibly give? Surely, if I have rightly understood your goodness and prudence, I am not far wrong, were you to be consulted, in what your reply would be: first, to wait; then, if he should seem to be in hiding and to be putting things off, to consult friends, to ask who is acting as agent, to give notice at the house. It is hard to say how many things you would all answer ought to be done before having recourse, of necessity, to that last extreme.

\ciceroSection{55} What does Naevius say to all this? He laughs at our madness, no doubt, in supposing that he should be measured by the rule of the highest duty, in expecting in his life the institutions of honourable men. ``What have I to do with that lofty sanctity and scrupulousness? Let honourable men look to those duties; for me, the question is what I have, not by what means I have got it, nor how I was born, nor how I was brought up.'' I remember the old saying: a buffoon can become rich much more easily than the head of a household.

\ciceroSection{56} What he does not dare to say in words, in deed he openly proclaims. For if he wishes to live by the rule of good men, there are many things he must learn and unlearn, both of which are hard at his age. ``I did not hesitate,'' he says, ``when the bond had been broken, to advertise the goods.'' Wickedly! But since you claim this for yourself and demand it be granted, let it be granted. What if he never broke it? If this whole charge has been concocted from beginning to end by your sheer fraud and malice? If you had no bond with Publius Quinctius at all? What name shall we call you by? Wicked? But even if the bond had been broken, in this petition and advertisement of the goods you would still have been found utterly wicked. Malicious? You do not deny it. Fraudulent? You actually claim it for yourself and think it praiseworthy. Bold, greedy, perfidious? Common and stale words; the matter is new and unheard-of.

\ciceroSection{57} What is to be said, then? I am afraid, by Hercules, of using either weightier words than the natural force of the case bears, or lighter words than the case demands. You say the bond was broken. Quinctius asked you, the moment he returned to Rome, on what day you said this bond had been due. You replied at once: ``On the Nones of February.'' As Quinctius was leaving you, he came to remember the day on which he had set out from Rome for Gaul. He turns to his diary; he finds the day of his departure: the day before the Kalends of February. If he was at Rome on the Nones of February, we say nothing as to whether or not he gave you a bond.

\ciceroSection{58} What now? How can this be discovered? Lucius Albius set out with him --- a man of the first rank for honour. He will give his testimony. Friends accompanied both Albius and Quinctius. They too will give their testimony. The diary of Publius Quinctius, the witnesses --- so many, with the strongest possible reasons why they could know, none why they should lie --- will be set against your single corroborator.

\ciceroSection{59} And in a case like this, will Publius Quinctius struggle? Will the wretched man go on longer in such fear and danger? Will his adversary's standing terrify him more vigorously than the judge's good faith consoles him? For he has always lived a rough and uncultivated life. He has been by nature austere and reserved. He has not been one for the sundial, nor for the Campus, nor for dinner-parties; his aim has been to keep his friends by attention, his estate by parsimony. He has loved the old idea of duty, the splendour of which has been worn dim by these manners of ours. If, in an equal case, he were to seem to leave the field the worse off, even then there would be no slight cause for complaint. As it is, in the better case he does not even ask to be on equal terms; he allows himself to be on inferior terms --- only just so far as not to be surrendered, with all his goods, his reputation, his fortunes, to the appetite and cruelty of Sextus Naevius.

\ciceroSection{60} I have shown what I promised first, Gaius Aquilius: that there was no cause whatsoever for Naevius to make his application; because nothing was owed in money, and that, even granting it had been owed, no act had been done that would warrant resort to that procedure. Now turn your attention to the point that Publius Quinctius's goods could in no way be possessed under the praetor's edict. Read the edict. ``The man who shall have hidden himself for the sake of fraud.'' Quinctius is not such a man --- unless those count as hiding who go about their own business and leave an agent in their place. ``The man who shall leave no heir.'' Not he, either. ``The man who shall have changed his country for the sake of exile.'' At what time, Naevius, do you think Quinctius ought to have been defended in his absence, and how? At the time when you were applying to take possession of the goods? Nobody was present, for nobody could have divined that you were going to apply, nor could anyone have been concerned to oppose what the praetor was ordering --- not, indeed, ordering to be done, but ordering to be done under his edict.

\ciceroSection{61} What was the first opportunity, then, given to the agent for defending the absent man? When you were posting the advertisements. Well: he was there, he did not allow it; Sextus Alfenus tore down the bills. The first step of duty was kept, and kept by the agent with the highest diligence. Let us see what came next. You catch a slave of Publius Quinctius's in public; you try to take him off. Alfenus does not allow it; he wrests the man from you by force; he sees to it that the slave is brought back to Quinctius's house. Here too the diligent agent's duty stands fully maintained. You say Quinctius is in your debt. The agent denies it. You wish him to give bond; he gives it. You summon him to court; he follows. You demand a trial; he does not refuse. What it is, in any other terms, to be defended in absence, I do not understand. ``But what sort of person was the agent?''

\ciceroSection{62} Some shabby fellow picked out, no doubt, a litigious wretch, capable of bearing a wealthy buffoon's daily abuse. Nothing of the kind. A wealthy Roman knight, well in command of his own affairs --- in fact the very man whom, every time Naevius set out for Gaul, he left as his own agent at Rome. And do you dare, Sextus Naevius, to deny that Quinctius was defended in his absence, when the very man who used to defend you was defending him? When the man into whose hands you used, on setting out, to commend and entrust your property and your name accepted the trial on Quinctius's behalf, do you try to say that no one was found to defend Quinctius in court?

\ciceroSection{63} ``I demanded,'' he says, ``that he give security.'' You demanded it wrongly. ``It was so ordered.'' Alfenus refused. ``So, but it was the praetor who decreed.'' --- The tribunes were therefore appealed to. --- ``At this point I have you,'' he says, ``this is not to suffer trial nor to defend by trial, when you seek aid from the tribunes.'' This, when I consider how prudent Hortensius is, I do not think he will say. But when I hear that he has said it before, and consider the case itself, I cannot find what else he can say. For he admits that Alfenus tore down the bills, gave bond, did not refuse to accept the trial in the very words Naevius proposed --- only by the established custom and procedure, through that magistrate who is appointed for the very purpose of giving aid.

\ciceroSection{64} Either these things must not have been done, or else Gaius Aquilius, a man of his quality, must, on his oath, lay down this rule of right in our state: that one whose agent has not accepted every trial put up to him in any words a plaintiff likes; one whose agent has dared to appeal from the praetor to the tribunes --- that man is not defended; his goods can rightly be possessed; and it is fitting that, in his absence, miserable, ignorant of his fortunes, every ornament of his life should be torn from him with the highest disgrace and ignominy.

\ciceroSection{65} Since this can be approved by no one, this surely must be approved by everyone: that Quinctius, in his absence, was defended in court. And that being so, the goods were not possessed under the edict. ``But the tribunes did not even hear the case.'' I admit that, if so, the agent ought to have obeyed the praetor's decree. What of it, then? If Marcus Brutus said openly that he would interpose his veto, unless something were settled between Alfenus himself and Naevius --- does it not appear that the appeal to the tribunes succeeded, not for delay, but for assistance?

\ciceroSection{66} What happens next? Alfenus --- so that all might understand that Quinctius was being defended in court, and lest any base suspicion could attach either to his own duty or to Quinctius's reputation --- summons several honourable men, attests in Naevius's hearing that on the basis of their common kinship he asks first that Naevius do nothing more savage in his absence to Publius Quinctius without cause; but if, on the other hand, he is determined to press the contest with all hostility, that Alfenus is ready, by every honest and rightful procedure, to defend the proposition that what he claims is not owed; he accepts the trial on the terms Naevius proposes.

\ciceroSection{67} The tablets of that proposal and condition were sealed by several honourable men. The matter cannot fall into doubt. So, with everything still entire, with the goods neither advertised nor in possession, Alfenus undertakes that Quinctius will appear before Naevius. Quinctius comes to that bond. The case lies in dispute, with Naevius wrangling, for two whole years, until a way could be found by which the matter might depart from the usual practice and the whole cause be locked into this unique trial.

\ciceroSection{68} What duty of an agent, Gaius Aquilius, can be mentioned that Alfenus appears to have neglected? On what grounds is it claimed that Publius Quinctius was not defended in absence? Or perhaps Hortensius will say --- as I think he will, since he has lately suggested it and Naevius keeps shouting it --- that Naevius did not have an equal contest with Alfenus at that time, when those persons were dominant. Even if I should be willing to admit it, they will surely concede this much: it was not that Publius Quinctius had no agent, but that his agent had influence. For my purposes, it is enough for victory that he had \emph{an} agent through whom the trial could be brought. What sort of man he was, provided he was defending the absent man through law and through a magistrate, has nothing to do with the matter, in my judgment.

\ciceroSection{69} ``He belonged to that party,'' he says. Why not? --- a man brought up in your house; a man you had so trained from boyhood that he would not back even a noble gladiator. If Alfenus wanted what you yourself always wanted to the highest degree, were you on that score not on equal terms with him? ``He was a friend of Brutus's,'' he says; ``and so Brutus was the one interposing.'' But you, on the other side, had Burrienus deciding wrong, and indeed all those who at that time had the greatest power through force and crime, and who dared what they could. Or did you wish to defeat all those whom these now labour with so much effort to make you defeat? Have the courage to say so --- not in public, but to the very men you have summoned.

\ciceroSection{70} Yet I am unwilling to revive that subject by recalling it; the memory of which I think ought to be utterly torn up and erased. I will say only this: if Alfenus was strong because of party loyalty, Naevius was the strongest. If Alfenus, relying on influence, was demanding something inequitable, Naevius was getting much more inequitable things. There was no difference, I think, between your party allegiance: in talent, seniority, technique, you easily defeated him. To leave the rest aside, this is enough: Alfenus perished along with and on account of those whom he loved; you, on the other hand, when those who had been your friends could no longer prevail, contrived that those who did prevail should be your friends.

\ciceroSection{71} But if you do not think you had equal right at that time with Alfenus --- since, however, he could call out someone against you, since some magistrate could be found before whom Alfenus's case could stand --- what must we say at the present time of Quinctius? For whom no fair magistrate has yet been found, no usual trial granted; no terms, no sponsio, no demand has ever been made --- and not just no fair demand, but a demand of a kind that, until this time, no one had even heard of. ``I want to dispute about money'' --- it is not allowed. ``But that is the controversy.'' --- ``It does not concern me; you must plead a capital cause.'' --- ``Make your accusation, then, where that is necessary.'' --- ``No,'' he says, ``unless you first speak in the prior place, in this new way.'' --- ``Then I must perforce speak.'' --- ``Hours will be set in advance at our pleasure; the judge himself will be coerced. ---''

\ciceroSection{72} ``What then?'' --- ``You will find some advocate or other, a man of old-school duty, who would not give a thought to our brilliance and influence. Lucius Philippus will fight for me --- the most flourishing man in our state for eloquence, weight, and honours; Hortensius will speak --- a man of outstanding talent, nobility, and reputation. Beside, the noblest and most powerful men will be present: their numbers and their massed presence will make even him quail who is a stranger to the danger, let alone Publius Quinctius, whose head is at stake.''

\ciceroSection{73} This is an unequal contest, not the one in which you used to ride against Alfenus. To this man you have not even left a place to stand against you. So either you must show that Alfenus denied he was the agent, did not tear down the bills, was unwilling to accept the trial; or, since these things were so done, you must concede that you did not take possession of Publius Quinctius's goods under the edict. For if you had taken possession under the edict, I ask: why were the goods not sold? Why did the other guarantors and creditors not turn up? Was no one to whom Quinctius owed anything? There were, and there were a great many, because his brother Gaius had left a fair amount of debt. What of it, then? They were all complete strangers to our man, and they were owed money; but no one was found so signally wicked as to dare violate the reputation of Publius Quinctius in his absence.

\ciceroSection{74} One man only there was --- a kinsman by marriage, a partner, a connection, Sextus Naevius --- who, though he himself was a debtor besides, contended as eagerly as if a special prize had been offered for the crime, that the kinsman whom he had himself afflicted and overthrown should be deprived not only of his honourably acquired goods but even of the common light of day. Where were the other creditors then? Where, indeed, are they now? Is anyone saying that he hid himself for the purpose of fraud? Anyone denying that Quinctius in his absence was defended? No one is to be found.

\ciceroSection{75} On the contrary, all those who have, or have had, dealings with him are present, supporting him, working that the loyalty of our man --- known on so many occasions --- be not undermined by the perfidy of Sextus Naevius. For a sponsio of this kind, witnesses ought to be supplied from the kind of men who would say: ``He broke a bond with me; he defrauded me; he asked me for a postponement on a debt he had denied; I could not bring him to a trial; he hid himself; he left no agent.'' None of these things is being said. Witnesses are being procured to say so. Well, we will see, when they have spoken. Yet let them ponder this one thing: that they are weighty enough to be able, if they choose to keep the truth, to keep their weight; if they neglect it, they are slight enough that all may understand their authority is no help in maintaining a lie, only in proving the truth.

\ciceroSection{76} I ask these two things: first, on what reasoning Naevius did not finish the business he had undertaken --- that is, why he did not sell the goods of which he was in possession under the edict; and second, why, of so many creditors, no one else came forward on this account. From which it must necessarily be admitted either that no one else was so reckless, or that you yourself, after most shamefully undertaking the thing, could not bring yourself to persevere and complete it. What if you, Sextus Naevius, have yourself decided that Publius Quinctius's goods were not possessed under the edict? Your own testimony, which would be worthless in another's case, ought to weigh most heavily in your own, since it tells against you. You bought the goods of Sextus Alfenus when Lucius Sulla, the dictator, was selling them; you registered Quinctius as your partner in the purchase. I say no more. Were you forming a voluntary partnership with the man who, you say, had defrauded you in your inherited partnership? Were you commending --- by your own judgment --- the man whom you supposed stripped of all reputation and fortune?

\ciceroSection{77} I doubted, by Hercules, Gaius Aquilius, whether I could stand fast in this case with steady and confirmed mind. So I thought to myself: when Hortensius is going to speak against me, when Philippus is going to listen to me attentively, I shall, in many points, slip from fear. I told this Quintus Roscius, whose sister is married to Publius Quinctius, when he was asking and most earnestly pressing me to defend his kinsman, that it was very hard for me, against such orators as these, not merely to make my peroration in such a case but even to attempt to utter a word at all. Since he was urging eagerly, I told him --- as friend to friend --- that anyone with a brazen face who tries to do gestures with him in the audience seems to me to be impudent, and that anyone who contends with him --- even if before they had seemed to have something correct or graceful --- loses it. Lest something of the same kind happen to me, when about to speak against such an artist, I am afraid.

\ciceroSection{78} Then Roscius told me many other things to encourage me; so that, by Hercules, even if he had said nothing, by his own silent duty and zeal toward his kinsman he would have moved anyone --- for as he is an artist of a kind that he alone seems worthy to be looked at on stage, so he is a man of a kind that he alone seems worthy not to come on it. Yet at last: ``What,'' he said, ``if you have a case in which you must make this plain, that no one can in two days, or at the outside three, walk seven hundred miles --- are you still afraid you cannot make this contention against Hortensius?''

\ciceroSection{79} ``Not at all,'' I said; ``but what has that to do with the matter?'' ``Beyond doubt,'' he said, ``the case rests on it.'' How? He is teaching me one of those matters of fact and acts of Sextus Naevius which, if it were brought forward together, would by itself be enough. I beg you, Gaius Aquilius, and you who are sitting in council, to attend carefully. You will surely understand that on the one side, from the start, appetite and audacity have been fighting; on the other, truth and modesty have been holding out as far as they could. You apply for leave to take possession of the goods under the edict. On what day? You yourself, Naevius, I want to hear; I want a deed unheard-of to be convicted by the very voice of the man who committed it. Tell us the day, Naevius. ``On the fifth day before the Kalends of the intercalary month.'' Well said. How far from here is it to your Gallic estate? Naevius, I am asking you. ``Seven hundred miles.'' Excellent. Quinctius is thrown off the estate --- on what day? Can we hear this too from you? Why are you silent? Tell us the day, I say. He is ashamed to say. I see. But it is too late and useless to be ashamed. He is thrown off the estate, Gaius Aquilius, the day before the Kalends of the intercalary month. Two days later, or, if someone ran straight from the law-court at the very moment, in not three full days, seven hundred miles cannot be covered.

\ciceroSection{80} O incredible thing! O reckless appetite! O winged messenger! Sextus Naevius's footmen and satellites cross from Rome over the Alps into the Sebagnine country in two days. O lucky man, to have such messengers --- no, rather, such Pegasi! Here, even if all the Crassi and Antonii of the world should rise up, even if you, Lucius Philippus, who flourished among them, were willing to plead this case with Hortensius, still I must come off the better. For things are not, as you suppose, all in eloquence. There is a kind of truth so visibly clear that nothing can shake it.

\ciceroSection{81} Did you, before you applied to take possession of the goods, send to see to it that the owner be forcibly thrown out of his own estate by his own household? Choose either: one is incredible, the other a crime, and until now both unheard-of. Do you maintain that seven hundred miles were covered in two days? Speak. You say no? Then you sent before. So much the worse. If you said the first, you would seem to be lying impudently; in admitting the second, you grant that you have committed a thing that not even a lie can cover up. Will so eager, so audacious, so reckless a piece of policy be approved by Aquilius and by men of his standing?

\ciceroSection{82} What does this madness, this haste, this immaturity signify, but force, but crime, but brigandage --- in short, anything sooner than law, than duty, than modesty? You sent without the praetor's order. With what design? You knew he was going to order it. What of it? Once he had ordered, could you not then send? You were going to apply. When? After thirty days. Of course, that is, if nothing hindered you, if your wish remained the same, if you were in good health, if you were even alive at all. Of course, the praetor would have given the order --- yes, if he willed, if he were in health, if he were on the bench, if no one objected, if the man would, by his own decree, give security and accept trial.

\ciceroSection{83} For, by the immortal gods! if Alfenus, Publius Quinctius's agent, had at that moment been willing to give you security and accept the trial, willing in fact to do everything you were applying for, what would you have done? Would you have called back the man you had sent into Gaul? But by then the man had already been thrown off his estate, hurled out headlong from his own household gods, and --- what is most outrageous --- violently assaulted by the hands of his own slaves, on your message and at your bidding. You were going to put these things right afterwards, of course. Do you dare to speak of any man's life --- you, who must concede that you were so blinded by appetite and avarice that, ignorant of what would later happen, with many things capable of intervening, you set the hope of an immediate wrong on the uncertain outcome of the time to come? And I am speaking as if at the very moment when the praetor had ordered you to take possession under the edict, if you had sent men into possession, you ought to have had, or could have had, the right of throwing Publius Quinctius out of possession.

\ciceroSection{84} Everything in this case, Gaius Aquilius, is so plain that anyone can see in it wickedness and influence pitted against want and truth. How did the praetor order you to take possession? Under his edict, of course. In what words was the sponsio framed? ``If Publius Quinctius's goods were not in possession under the praetor's edict.'' Let us go back to the edict. How does it order possession to be taken? Is there any reason, Gaius Aquilius, why, if Naevius took possession in a way far different from what the praetor's edict prescribed, he did not take possession under the edict, and I have won the sponsio? None, I think. Let us examine the edict. ``Those who under my edict shall come into possession.'' He is speaking of you, Naevius, in your own way of looking at it; for you say you came under the edict. He is defining what you are to do, instructing you, giving you precepts. ``It seems fitting that they be in possession in this manner.'' How? ``What they will rightly be able to take care of in the place itself, let them take care of in the place itself; what they cannot, that it shall be lawful for them to carry away and lead off.'' What then? ``It is not our pleasure that the unwilling owner be thrust out.'' Even the very man who is hiding himself for fraud, even the very man whom no one has defended in trial, even the very man who is dealing badly with all his creditors --- it forbids him to be thrust out unwilling from the property.

\ciceroSection{85} As you set out to take possession, the praetor himself, Sextus Naevius, says to you in plain terms: ``See to it that you possess in such a way that Quinctius possess together with you; possess in such a way that no force is brought against Quinctius.'' What of it? How do you observe this? I forgo the point that the man was not hiding himself, that he had a house at Rome, a wife, children, an agent; the man who had not failed to come to your bond. All this I forgo. I say only this: that the owner was driven out of his estate; that the owner had hands laid on him by his own household before his own household gods. I say this: if anyone takes possession by some means or other of one estate, but allows the owner to keep the rest of his property, he seems --- in my judgment --- to be in possession of the estate, not of another man's goods. What is it to take possession? Surely it is to be in possession of the things that, at the time, can be possessed. When there was a house at Rome, slaves, and in Gaul itself private estates of Publius Quinctius which you never dared to take into possession; if you were taking possession of Publius Quinctius's goods, you ought to have been in possession of them all under the same right. (\textit{From a paraphrase by Iulius Severianus, ch. 16:} that Naevius did not even address Quinctius, when he was both at hand and able to bring suit daily; secondly, that he preferred all those most difficult trials, brought with the greatest unpopularity to himself and the greatest danger to Publius Quinctius, to the one suit about money that could be settled in a single day --- from which, by his own admission, all this has sprung.) On which point I made the offer: that if he wished to claim money, Publius Quinctius would give security for satisfying judgment, provided that he himself, if he wanted to claim anything, used the same condition.

\ciceroSection{86} I have shown how many things ought to have happened before the goods of a kinsman could be applied for in possession --- especially when his house, wife, and children were at Rome and an agent who was kinsman to both. I have shown that, although Naevius says the bond was broken, there was no bond at all: that on the day on which he says Quinctius gave it to him, Quinctius was not even at Rome. I have promised to make this plain by witnesses who would have every reason to know and no reason to lie. I have shown, further, that the goods could not have been possessed under the edict, since Quinctius had neither hidden himself for the sake of fraud nor was said to have changed his country for the sake of exile.

\ciceroSection{87} It remains: that no one defended him at trial. To which I have replied that on the contrary he was abundantly defended --- not by some stranger, nor by some pettifogger or rogue, but by a Roman knight, his own kinsman and connection, the very man Sextus Naevius himself had been accustomed before now to leave as his agent. Nor was he, by appealing to the tribunes, any the less prepared to suffer a trial; nor was Naevius's right wrenched from him by the agent's influence. On the contrary, Naevius was at that time only just superior in his own influence, and now scarcely allows us breathing-room.

\ciceroSection{88} I have asked what cause there was that the goods were not sold, when they were in possession under the edict; and further, why, of so many creditors, none did the same thing then nor speaks against him now, but all are fighting for Publius Quinctius --- particularly when in such a trial the testimony of creditors is reckoned to bear most on the matter. Then I have used the testimony of the adversary himself, who recently registered as his partner the very man whom now, on his pleading, he claims was not even among the living. Then I produced that incredible speed --- or rather audacity. I established the necessity that either seven hundred miles were covered in two days, or that Sextus Naevius sent men into possession many days before he applied for leave to take possession of the goods.

\ciceroSection{89} I then read the edict, which expressly forbids that an owner be thrust off an estate. By which it is settled that Naevius did not take possession under the edict, since he admits that Quinctius was forcibly thrust off the property. I established further that the goods were not, on the whole, possessed --- because the possession of goods is reckoned not in some part but in the whole of what can be held and possessed. I have said that there was a house at Rome which Naevius did not even approach; many slaves, of whom he took possession of none, did not even touch one. There was one whom he tried to touch; he was prevented and let it go.

\ciceroSection{90} You learnt that in Gaul itself Sextus Naevius did not enter the private estates of Quinctius; lastly, that even from this estate, of which he took possession after forcibly expelling his partner, not all of Quinctius's private slaves were ejected. From this, and from the rest of Sextus Naevius's words, deeds, and designs, anyone can understand that he has done nothing else, and is doing nothing else now, except contrive that by violence, by injustice, by inequity of judgment he can make his own the whole estate that is held in common.

\ciceroSection{91} Now that the case has been pleaded out, the matter itself and the magnitude of the danger seem to compel Publius Quinctius, Gaius Aquilius, to beseech and entreat you and those who sit beside you in your council --- by his old age, by his isolation --- to do nothing other than to follow the bent of your own nature and goodness; that, since the truth makes with him, his want should weigh more towards mercy than this man's wealth towards cruelty.

\ciceroSection{92} On the day we came to you as judge, on that same day we began to set lightly by those threats which we had previously dreaded. If case were contending with case, we judged that we could readily prove ours to anyone. The reason it was, that one rule of life was to be set against another rule of life is the very reason we judged we needed all the more to have you for judge. For the question is now in the balance: whether that rough and uncultivated parsimony of his can defend itself against luxury and licence, or whether, defaced and stripped of every ornament, it must be surrendered naked to appetite and insolence.

\ciceroSection{93} Publius Quinctius does not match himself with you in influence, Sextus Naevius; he does not contend in resources, in capacity. All the arts in which you are great he yields to you; he confesses that he cannot speak elegantly, cannot speak to please, cannot desert a friendship in trouble and fly off to one in flower; cannot live by lavish expense; cannot adorn a banquet with magnificent splendour; cannot have a house closed to modesty and sanctity but open --- in fact thrown open --- to appetite and pleasures. He claims, on the other hand, that duty, faith, diligence, a life always rough and dry, have been dear to his heart. That those things are higher today and count for most in our manners --- he is aware. What of it, then?

\ciceroSection{94} Not so far that those who, abandoning the discipline of good men, have preferred to follow Gallonius's manner of profit and expenditure --- and who, what was not in him, have moreover lived with audacity and perfidy --- should be lord over the lives and fortunes of the most honourable men. If a man may live whom Sextus Naevius does not wish to, if there is room for an honourable man in the state against Naevius's will, if it is right for Publius Quinctius to draw breath against Naevius's nod and command, if the marks of distinction his modesty has earned for him he can keep, against insolence, with me defending him, then there is hope that this miserable and unhappy man may at last be allowed to stand. But if Naevius shall be permitted to do whatever he likes, and shall like what is not lawful, what is to be done? What god is to be called on? What man's faith is to be implored? What complaint, what mourning, what sorrow can fitly be found, in so great a calamity?

\ciceroSection{95} It is miserable to be tipped off all one's fortunes; more miserable, by injustice. It is bitter to be cheated by anyone; more bitter, by a kinsman. It is calamitous to be overturned in one's goods; more calamitous, with disgrace. It is deadly to have one's throat cut by a brave and honourable man; more deadly, by one whose voice was once for sale on the auctioneer's stand. It is unworthy to be defeated by an equal or a superior; more unworthy, by an inferior or a baser man. It is grievous to be handed over to another along with one's goods; more grievous, to an enemy. It is horrible to plead a capital case; more horrible, to plead it in the prior place.

\ciceroSection{96} Quinctius has looked all about, has tried everything, Gaius Aquilius. He could not find a praetor from whom he could obtain his right --- not even one before whom he could lodge his application as he pleased. He could not even find friends of Sextus Naevius --- at whose feet he often and at length lay prostrate, beseeching them by the immortal gods either to contend with him by law or to lay any injustice on him without ignominy.

\ciceroSection{97} At last he has subjected himself to that adversary's most haughty face. He has, with tears, taken Sextus Naevius's own hand --- the hand practised in advertising a kinsman's goods --- and besought him by the ashes of his own dead brother, by the name of kinship, by Naevius's own wife and children (none of whom is closer than Publius Quinctius), that at last he show some mercy: have some regard, if not for kinship, then for his own age; if not for the man, then for human nature. That he settle with him on any terms whatever, with his reputation kept whole --- on any terms, however hard, only let them be tolerable.

\ciceroSection{98} Refused by Naevius himself, unsupported by Naevius's friends, harassed and terrified by every magistrate, he has no one whom he can appeal to except you. To you he commits himself, to you all his goods and fortunes; to you he commits the reputation and hope of the rest of his life. Worn down by many insults, tossed by many wrongs, he takes refuge with you not as a disgraced man but as a wretched one. Thrown out of his richest estate, set upon by every kind of disgrace, when he saw Naevius lording it in the goods his own father bequeathed, he could not provide a dowry for his marriageable daughter --- yet he committed nothing alien to his earlier life.

\ciceroSection{99} Therefore he beseeches this of you, Gaius Aquilius: that the reputation, the honour which now, when his life is almost over and run out, he has brought into your judgment, he be allowed to bear away with him from this place --- that the man on whose duty no one ever cast doubt should not, in his sixtieth year, be marked with disgrace, with a stain, with the foulest ignominy; that Sextus Naevius should not abuse all his ornaments as if they were spoils; that it should not happen through you that the same reputation which carried Publius Quinctius all the way to old age should not also escort him all the way to the funeral pyre.

